% Contrepartie francaise de paper-cleveref.tex. Elle garde honnetes deux
% choses que rien d'autre ne surveille :
%
%   1. babel-french rend « : » actif, et toutes les etiquettes du template en
%      contiennent un (def:, thm:, prop:, cor:). Sans le garde-fou de
%      ocots-env.sty, cleveref n'arrive plus a les relire : « Missing
%      \endcsname inserted » en rafale, et des renvois rendus
%      « :admissible@cref ». Voir l'issue #44.
%   2. la table des chaines : les intitules de renvoi viennent de
%      \ocotsstring, donc une traduction manquante se verrait ici, et nulle
%      part ailleurs (check-crefnames.sh detecte le marqueur « ?clef? »).
%
% Noter que cleveref est charge NU, sans option de langue : c'est le fichier
% lang/ qui la lui passe par \PassOptionsToPackage. C'est ce qui fait que la
% conjonction est « et » et non « and » ci-dessous -- comportement teste ici.
%
% Renvois cleveref : chaque famille d'enonce doit se citer sous SON intitule.
%
% Les cinq familles de boites (theorem, definition, proposition, corollary,
% conjecture) partagent le compteur `theorem'. Rien dans « make check » ne
% pouvait voir qu'elles se citaient toutes comme « Theorem » : le document
% compile sans le moindre avertissement. D'ou ce fichier et son garde-fou
% check-crefnames.sh -- meme dispositif que check-numbering.sh pour le
% partage de compteur.
%
% Convention lue par check-crefnames.sh : une ligne
%   CREFCHECK <Intitule attendu> = \cref{...}
% doit se rendre en « CREFCHECK <Intitule attendu> = <Intitule attendu> N.M ».
\documentclass[11pt]{ocots-paper}
\usepackage[lang=fr]{ocots}
\usepackage{cleveref}

\newcommand{\R}{\mathbb{R}}
\title[Exemple cleveref]{Renvois cleveref et familles d'\'enonc\'es}
\author{Ada Auteur}
\address{D\'epartement de math\'ematiques, Universit\'e Exemple}
\email{ada@example.org}

\begin{document}
\maketitle

\section{\'Enonc\'es}

\begin{definition}[Admissible point]
\label{def:admissible}
A point is admissible when it is smooth on $\R$.
\end{definition}

\begin{proposition}[Stability]
\label{prop:stability}
Admissible points form an open set.
\end{proposition}

\begin{lemma}[Compactness]
\label{lem:compactness}
Every admissible sequence has a convergent subsequence.
\end{lemma}

\begin{lemma}[Closedness]
\label{lem:closedness}
The limit of an admissible sequence is admissible.
\end{lemma}

\begin{theorem}[Main result]
\label{thm:main}
Every admissible point satisfies the conclusion.
\end{theorem}

\begin{corollary}[Consequence]
\label{cor:consequence}
The conclusion holds almost everywhere.
\end{corollary}

\begin{conjecture}[Open problem]
\label{conj:open}
The converse holds as well.
\end{conjecture}

\begin{hypothesis}[Regularity]
\label{hyp:regularity}
The data are smooth on $\R$.
\end{hypothesis}

\begin{hypothesis}[Boundedness]
\label{hyp:bounded}
The data are bounded on $\R$.
\end{hypothesis}

\section{Renvois}

\noindent CREFCHECK Définition = \cref{def:admissible}\par
\noindent CREFCHECK Proposition = \cref{prop:stability}\par
\noindent CREFCHECK Lemme = \cref{lem:compactness}\par
\noindent CREFCHECK Théorème = \cref{thm:main}\par
\noindent CREFCHECK Corollaire = \cref{cor:consequence}\par
\noindent CREFCHECK Conjecture = \cref{conj:open}\par
\noindent CREFCHECK Hypothèse = \cref{hyp:regularity}\par
\noindent CREFCHECK Lemmes = \cref{lem:compactness,lem:closedness}\par
\noindent CREFCHECK Hypothèses = \cref{hyp:regularity,hyp:bounded}\par

\end{document}
